\section{Introduction}

Customer churn---the rate at which subscribers discontinue their telecom service---remains one
of the most pressing challenges in the telecommunications industry. With monthly churn rates
typically ranging from 1.5\% to 3\%, a mid-sized carrier with 10 million subscribers can lose
150,000 to 300,000 customers each month. Given that acquiring a new customer costs five times
more than retaining an existing one, even modest improvements in retention yield substantial
profit gains. Research by Bain \& Company demonstrates that a 5\% improvement in retention
rates can increase profits by 25--95\%.

Traditional machine learning approaches to churn prediction have proven effective at
\textit{identifying} at-risk customers. Logistic Regression, Random Forests, and gradient-boosted
models can classify customers with reasonable accuracy based on behavioral and demographic
features. However, prediction alone is insufficient for two critical reasons. First, a
probability score does not explain \textit{why} a customer is at risk, leaving customer service
representatives without actionable context. Second, prediction does not prescribe \textit{what
to do}---the specific retention actions most likely to succeed for a given customer's profile.

This project addresses both gaps by layering a \textbf{Retrieval-Augmented Generation (RAG)}
pipeline and a \textbf{LangGraph agent} on top of an existing ML churn predictor. The RAG
component retrieves relevant retention strategies from a curated knowledge base of telecom
retention best practices. The LangGraph agent orchestrates a five-node state machine that
produces: (1) a churn probability and risk level from the ML model, (2) top risk factors
identified by deterministic rules, (3) retrieved context from the knowledge base, (4) a
plain-English explanation generated by an LLM, (5) three personalized, actionable retention
recommendations, and (6) a one-sentence executive summary.

The system operates on the IBM Telco Customer Churn dataset containing 7,043 customers
across 20 features including tenure, contract type, monthly charges, internet service type,
payment method, and various add-on services. The dataset has a churn rate of approximately
26.5\%, making it a moderately imbalanced binary classification problem. Five models were
benchmarked: Linear Regression (as a baseline), Logistic Regression, Decision Tree, Random
Forest, and XGBoost. Logistic Regression was selected as the production model based on the
highest F1-score of 0.609.

The Gradio web application provides two tabs: a Quick Predict tab that runs the ML model
instantly without any API key, and an AI Agent Analysis tab that invokes the full agent
pipeline through Groq's \texttt{llama3-70b-8192} LLM. This design ensures the core ML
functionality remains fully operational offline while the AI enhancements are available
when an API key is configured.
